\section{Where this goes}
\label{sec:short:where}

\subsection{The dates}

\begin{figure}[H]
\centering
\begin{tikzpicture}[font=\scriptsize,
  ev/.style={draw=FluxBlue, fill=FluxBlue!9, rounded corners=2pt, text width=25mm,
             align=center, inner sep=4pt},
  now/.style={draw=FluxSlate, fill=FluxSlate!10, rounded corners=2pt, text width=25mm,
              align=center, inner sep=4pt},
  tick/.style={FluxSlate, thick}, lead/.style={FluxSlate!55, thin}]
%% Axis fits the text block; events alternate between two heights so that
%% neighbouring boxes cannot collide however long their labels run.
\draw[->, thick, FluxSlate] (0,0) -- (13.2,0);
\foreach \x/\l in {0.7/{Sep 2026\\now}, 2.0/{Oct 2026}, 4.9/{Mar 2027}, 7.2/{Apr 2027}, 10.1/{Jul 2027}, 12.3/{Aug 2027}}
  {\draw[tick] (\x,-0.12) -- (\x,0.12);
   \node[below=2pt, text=FluxSlate, align=center] at (\x,-0.12) {\l};}
%% low row
\draw[lead] (0.7,0.12) -- (0.7,0.75);
\node[now, anchor=south] at (0.7,0.75) {measured baseline\\ \num{6489} nodes, \num{1195} apps};
\draw[lead] (7.2,0.12) -- (7.2,0.75);
\node[ev,  anchor=south] at (7.2,0.75) {shielded pool sunset\\ block \num{3600000}};
\draw[lead] (10.1,0.12) -- (10.1,0.75);
\node[ev,  anchor=south] at (10.1,0.75) {\textbf{PNR activates}\\ parallel-asset\\ mining cut\\ block \num{3787502}};
%% high row
\draw[lead] (2.0,0.12) -- (2.0,2.35);
\node[ev,  anchor=south] at (2.0,2.35) {\textbf{spec v9 enforced}\\ block \num{3050000}\\ inference benchmark};
\draw[lead] (4.9,0.12) -- (4.9,2.35);
\node[ev,  anchor=south] at (4.9,2.35) {retiring chains one-way\\ no swaps in\\ block \num{3450000}};
\draw[lead] (12.3,0.12) -- (12.3,2.35);
\node[ev,  anchor=south] at (12.3,2.35) {retiring chains shut down\\ residue to the Foundation\\ block \num{3880000}};
\end{tikzpicture}
\caption{The dated schedule. Heights are consensus facts and calendar dates assume the thirty-second
block target holds. Everything to the right of ``now'' is planned.}
\label{fig:short:timeline}
\end{figure}

Three things are true about this schedule and worth stating together. The heights are fixed and
public, so anyone can check whether they were met. The order is deliberate: the specification
change lands roughly eight months before the economic change that depends on it, so the fields PNR
needs have time to prove themselves. And the whole sequence is falsifiable --- the full paper names,
for each milestone, the public observation that would show it had not happened.

\subsection{How to participate}

\textbf{Run a node.} Lock the collateral for a tier, run the node software on hardware that clears
that tier's benchmark, and the chain does the rest. Income today is issuance by rotation; from July 2027 it is issuance plus a share of application revenue. What you cannot do is earn more by hosting
more --- see \Cref{sec:short:limits} before deciding that is what you want.

\textbf{Deploy an application.} Write a specification, sign it, pay for it. There is no account and
no approval step. The application runs on several nodes until the paid period ends, and renewing it
is your act. If you want subscriptions and cards, a third-party service offers them on top; the
protocol does not.

\textbf{Build for machines.} If your tenant is a program --- an agent, a pipeline, a system that
provisions and releases capacity on its own --- \Flux{} is being built for exactly that caller. Give
it a key, fund the key with what you are willing to lose, and it needs nothing else from anyone.

\textbf{Read the source.} The chain, the node software and the tooling are public repositories, and
the full paper cites them by file and line for every claim about behaviour. Where it cites a private
repository, it says so.
